% ===========================================================================
% rechaos: A Machine-Checked Test-Case Shrinker
% OSDI-style two-column layout, Computer Modern (default LaTeX fonts).
% Build:  make        (uses Nix-provided texlive; see Makefile)
% ===========================================================================
\documentclass[10pt,twocolumn,letterpaper]{article}

% --- OSDI/USENIX-ish geometry, Computer Modern kept (no font package) -------
\usepackage[letterpaper,margin=1in,columnsep=0.33in]{geometry}
\usepackage[T1]{fontenc}
\usepackage{microtype}
\usepackage{amsmath,amssymb,amsthm}
\usepackage{mathtools}
\usepackage{booktabs}
\usepackage{listings}
\usepackage{xcolor}
\usepackage{url}
\usepackage[hidelinks]{hyperref}
\usepackage{enumitem}
\usepackage[numbers,sort&compress]{natbib}

\setlength{\parindent}{1em}
\setlength{\parskip}{0pt}

% Tighter section headings in the OSDI idiom.
\usepackage{titlesec}
\titleformat{\section}{\normalfont\large\bfseries}{\thesection}{0.6em}{}
\titleformat{\subsection}{\normalfont\normalsize\bfseries}{\thesubsection}{0.6em}{}
\titlespacing*{\section}{0pt}{1.1ex plus .2ex}{0.6ex}
\titlespacing*{\subsection}{0pt}{0.9ex plus .2ex}{0.4ex}

\newtheorem{theorem}{Theorem}
\newtheorem{definition}{Definition}

% Lightweight code style for Haskell / Lean snippets.
\lstdefinestyle{src}{
  basicstyle=\ttfamily\footnotesize,
  breaklines=true,
  columns=fullflexible,
  keepspaces=true,
  showstringspaces=false,
  xleftmargin=0.5em,
  aboveskip=0.6ex, belowskip=0.6ex,
}
\lstset{style=src}

\newcommand{\rc}{\textsc{rechaos}}
\newcommand{\tool}[1]{\texttt{#1}}

\title{\bf A Machine-Checked Test-Case Shrinker:\\
Verified Termination and Witness Preservation\\
for Chaos-Testing the Remote Execution API}

% Contributors listed alphabetically; no precedence.
\author{
{\rm b7r6} \qquad {\rm Claude Fable 5 (Anthropic)}\\
\rc\ project
}

\date{}

\begin{document}
\maketitle
\thispagestyle{empty}

% ===========================================================================
\begin{abstract}
% ===========================================================================
Test-case shrinking --- reducing a failing input to a small witness that still
triggers the failure --- is the feature that turns a fuzzer's haystack into an
actionable bug report. Yet every production shrinker we are aware of is itself
unverified software: a delta-debugging loop whose two load-bearing properties,
\emph{termination} and \emph{witness preservation}, are assumed rather than
proved. We present the shrinker of \rc, a black-box chaos and determinism tool
for the Bazel Remote Execution API (REAPI). The shrinker is a pure candidate /
observe state machine that reduces a failing fault \emph{timeline} by two
mechanisms: delta-debugging chunk deletion at halving sizes down to singletons
(yielding deletion-1-minimality for a deterministic predicate), and a
domain-specific \emph{intensity} axis that weakens individual faults along a
finite, decreasing measure. We prove in Lean~4, with no \texttt{sorry},
\texttt{admit}, or added axiom, that (i) the search terminates under a
well-founded $(\mathrm{length}, \mathrm{summed\text{-}intensity})$ lexicographic
measure and (ii) acceptance preserves the witness: the best reproducer advances
only on a confirmed reproduction, and an inconclusive (flaky, timed-out) verdict
is never mistaken for one. The Lean port is held faithful to the Haskell
production code by a reference-leads differential: corpora generated by running
the real implementation are re-checked by the Lean kernel via
\texttt{native\_decide}. A survey of the shrinking literature places this design
in context: production shrinkers typically treat termination and witness
preservation as engineering folklore rather than proved properties, and we argue
that the verification artifact --- not the reduction algorithm, which is ordinary
delta debugging --- is what is worth reporting. We analyze which of three
well-studied upgrades --- iterate-to-fixpoint, probabilistic query selection, and
choice-sequence reduction --- a high-assurance tool should adopt.
\end{abstract}

% ===========================================================================
\section{Introduction}\label{sec:introduction}
% ===========================================================================
Remote-execution and remote-cache servers sit on the critical path of every
build in large organizations. They are concurrent, content-addressed, and
exactly the kind of system whose interesting bugs --- torn reads served as
complete, dishonest commit sizes, queue entries that never get collected ---
surface only under adversarial timing and malformed-but-legal wire traffic.
\rc\ is a gateway that produces exactly that traffic against an \emph{unmodified}
REAPI endpoint~\cite{reapi}, asserts the invariants a correct CAS and scheduler
must hold, and --- when an invariant is violated --- deterministically replays,
compares, and \emph{shrinks} the failure to a minimal witness.

Shrinking is where a chaos tool earns its keep. A campaign that finds a
correctness violation under a thousand injected faults is useless if the
engineer who must fix it cannot tell which fault mattered. The shrinker's job is
to delete everything incidental and hand back the smallest timeline that still
reproduces the \emph{same} failure signature. This is delta debugging, and the
algorithm is forty years old. What is \emph{not} old, and what this paper is
about, is the observation that the shrinker is itself a piece of software on the
critical path of the tool's central claim --- ``here is the minimal cause'' ---
and that two of its properties are safety-critical to that claim:

\begin{itemize}[leftmargin=1.2em,itemsep=1pt,topsep=2pt]
\item \textbf{Termination.} The candidate search must not loop forever. A
shrinker that can cycle --- proposing a ``smaller'' candidate it has already
rejected --- is a liveness bug in a tool whose users run it unattended.
\item \textbf{Witness preservation.} Every reduction step the shrinker
\emph{accepts} must still reproduce the original failure. A shrinker that
accepts a candidate because a flaky run happened to fail, or because a timeout
was misread as the bug, produces a confident, minimal, and \emph{wrong} witness
--- the worst possible output, because it looks authoritative.
\end{itemize}

Every production shrinker we surveyed --- QuickCheck, Hedgehog, Hypothesis,
C-Reduce, the delta-debugging family --- treats both properties as folklore.
They are true, usually, for good engineering reasons, but they are not
\emph{proved}. For a tool whose entire value proposition is ``this witness is
the cause, and you can trust it,'' we think that gap is worth closing.

\paragraph{Contributions.}
\begin{enumerate}[leftmargin=1.4em,itemsep=1pt,topsep=2pt]
\item A pure, total shrinker (\S\ref{sec:design}) for fault timelines that
combines structural chunk deletion with a domain-specific fault-intensity axis,
implemented as a candidate / observe state machine with no \texttt{IO} and no
partial functions.
\item A Lean~4 formalization (\S\ref{sec:proofs}) proving termination under a
well-founded lexicographic measure and witness preservation under a conservative
acceptance rule, with a no-\texttt{sorry} / no-added-axiom policy enforced by CI.
\item A reference-leads differential methodology (\S\ref{sec:differential}) that
keeps the proved Lean model faithful to the executable Haskell by re-checking
reference-generated corpora inside the Lean kernel.
\item A survey and synthesis (\S\ref{sec:related},\,\S\ref{sec:discussion}) of
the shrinking literature, situating this design and analyzing three candidate
upgrades against the constraint that the verification artifact must survive.
\end{enumerate}

% ===========================================================================
\section{Background: the shrinking design space}\label{sec:background}
% ===========================================================================
Delta debugging~\cite{zeller1999yesterday,zeller2002ddmin} is the root of the
field. Given a failing input and an \emph{interestingness test} (an oracle that
says whether a candidate still fails), \texttt{ddmin} partitions the input and
greedily removes subsets, guaranteeing a \emph{1-minimal} result: no single
remaining element can be removed without losing the failure. Crucially,
1-minimality is a \emph{local} optimum. Global minimization is NP-complete, so a
1-minimal witness can still admit a strictly smaller 1-minimal
witness~\cite{zeller2002ddmin,vince2022iterating}. \texttt{ddmin} further assumes
the oracle is \emph{monotone} and \emph{consistent}; real oracles --- flaky
tests, timeouts, infrastructure errors --- violate this, and the algorithm's
guarantees degrade silently when they do.

Two branches extend the root. \emph{Structure-aware} reducers
(HDD~\cite{misherghi2006hdd}, Perses~\cite{sun2018perses},
hoisting~\cite{vince2021hoisting}, Vulcan~\cite{vulcan2023}) move reduction onto
a parse tree so that every candidate is syntactically valid, avoiding the wasted
oracle queries a flat reducer spends on unparseable intermediates.
\emph{Property-based-testing} shrinkers take a different tack: QuickCheck's
type-directed \texttt{shrink}~\cite{claessen2000quickcheck} is manual and
famously non-compositional under monadic bind; Hedgehog's integrated
shrinking~\cite{hedgehog} threads a rose tree of shrinks through the generator;
and Hypothesis's \emph{internal} shrinking~\cite{maciver2020hypothesis} reduces
the \emph{choice sequence} that drives the generator rather than the value,
making every shrink valid-by-construction and generator-agnostic. The last is
formally shortlex optimization: minimize first by length, then lexicographically,
over a well-founded total order on choice sequences.

The deepest framing is information-theoretic. Adaptive delta debugging is
adaptive combinatorial group testing~\cite{dorfman1943,aldridge2019grouptesting}:
each oracle query is a pooled test, positive iff the pool contains a
failure-essential element, and the group-testing literature supplies
algorithm-independent lower bounds on query count. Probabilistic Delta Debugging
(ProbDD)~\cite{wang2021probdd} operationalizes this as sequential Bayesian
experimental design, maintaining a keep-probability per element and choosing the
next subset to maximize expected information --- at the cost of abandoning the
1-minimality guarantee.

% ===========================================================================
\section{The \rc\ shrinker}\label{sec:design}
% ===========================================================================
\rc\ records a run as a \emph{timeline}: a list of \emph{decisions}, each an
observed wire event paired with the fault (if any) injected into it. The fault
algebra is small and semantic: \tool{delay} holds a message, \tool{abort} ends
an RPC with a gRPC status, \tool{dribble} paces a byte stream, \tool{truncate}
shortens a payload, and \tool{corrupt} rewrites leading bytes length-preservingly
so the server must re-hash to detect the tampering. A failing timeline is the
starting point; the shrinker's output is a small timeline that a user-supplied
checker confirms still reproduces the \emph{same} failure signature.

\subsection{Verdicts and the acceptance rule}
The checker returns one of three verdicts, and the three-way split is
load-bearing:
\begin{lstlisting}[language=Haskell]
data Verdict = Triggers       -- same failure signature
             | DoesNotTrigger  -- ran, did not reproduce
             | Unknown         -- flake / timeout / infra
\end{lstlisting}
Only \texttt{Triggers} --- a reproduction of the \emph{same} signature, not
merely a nonzero exit --- advances the search. \texttt{DoesNotTrigger} and
\texttt{Unknown} both discard the candidate. Treating \texttt{Unknown}
conservatively is exactly the defense against the witness-corruption failure mode
described in \S\ref{sec:introduction}: a flake never gets to masquerade as the
bug.

\subsection{The state machine}
The shrinker is a pure state machine over
$\mathit{ShrinkState} = (\mathit{best}, \mathit{pending})$, where $\mathit{best}$
is the smallest timeline confirmed to trigger and $\mathit{pending}$ is a queue
of candidates to try against it:
\begin{lstlisting}[language=Haskell]
start    :: Timeline -> ShrinkState
candidate:: ShrinkState -> Maybe Timeline
observe  :: Verdict -> ShrinkState -> ShrinkState
\end{lstlisting}
\texttt{start} seeds $\mathit{best}$ with the failing timeline's injected faults
and enqueues its candidates. \texttt{candidate} offers the head of the queue.
\texttt{observe Triggers} promotes the candidate to $\mathit{best}$ and
\emph{re-seeds} \texttt{pending} from the smaller timeline --- a greedy restart
that is what drives convergence; any other verdict pops the candidate and moves
on. The module imports no \texttt{IO}, has no partial indexing, and uses no
\texttt{error} or \texttt{undefined}.

\subsection{Candidate generation}
\texttt{candidates} is the union of two families:
\begin{itemize}[leftmargin=1.2em,itemsep=1pt,topsep=2pt]
\item \textbf{Deletions:} contiguous chunk removals at halving sizes
$n, n/2, \dots, 1$. Sweeping down to singletons is what makes a converged result
\emph{deletion-1-minimal} for a deterministic predicate --- no single retained
fault can be dropped without losing the failure.
\item \textbf{Intensity reductions:} replace one injected fault with a strictly
weaker variant via \texttt{weaker}. \tool{delay} halves its micros toward zero;
\tool{dribble} doubles its rate toward the full-speed cap
($\mathit{messageBytes} \times \mathit{microsPerSecond}$); \tool{truncate} raises
its kept-byte count toward the message size; \tool{corrupt} halves its corrupted
byte count toward one; \tool{abort} has no weaker form.
\end{itemize}
The intensity axis is the domain-specific part: a timeline can be structurally
minimal (every remaining fault necessary) yet still carry a 100\,ms delay where
1\,ms suffices, or a corruption of 64 bytes where 1 byte suffices. Intensity
reduction shrinks the \emph{magnitude} of the cause, not just its support.

% ===========================================================================
\section{Verification}\label{sec:proofs}
% ===========================================================================
The shrinker is ported to Lean~4~\cite{leanprover} as
\texttt{Rechaos.Minimize}, mirroring the Haskell definitions
constructor-for-constructor. Two abstract theorems are proved over core Lean (no
Mathlib, no \texttt{UInt64} arithmetic forced in the kernel), and the concrete
behavior is pinned to the reference by a differential corpus (\S\ref{sec:differential}).

\subsection{Termination}
Define the measure $\mu(t) = (|t|, \sigma(t))$ where $|t|$ is the timeline length
and $\sigma(t)$ is the summed intensity (severity) of its injected faults,
ordered lexicographically. Every generated candidate is strictly smaller than its
input under $\mu$:

\begin{theorem}[weaker\_severity\_lt]
For every fault $f'$ that \texttt{weaker}~$e~f$ emits, the severity of $f'$ at
anchoring event $e$ is strictly less than that of $f$.
\end{theorem}
The proof is by cases on the fault, discharging each branch's numeric
side-condition (e.g.\ for \tool{truncate}, the new kept count
$keep + \max(1, (\mathit{cap} - keep)/2)$ strictly increases yet stays within the
message, so the bytes-dropped severity strictly decreases). Together with
additivity of $\sigma$ over the \texttt{take ++ [replacement] ++ drop} splice
(\texttt{summedIntensity\_append}) this gives the intensity decrease; a nonempty
contiguous deletion strictly shrinks the length (\texttt{deletion\_length\_lt}).
Hence every candidate strictly descends $\mu$, which is well-founded, so the
shrink loop terminates.

\subsection{Witness preservation}
\begin{theorem}[observe\_best\_triggers]
If \texttt{observe}~$v~s$ changes $\mathit{best}$, then $v = \mathtt{Triggers}$
and the new $\mathit{best}$ is exactly the pending head the checker judged.
\end{theorem}
Contrapositively, $\mathit{best}$ advances \emph{only} on a confirmed
reproduction. Given a truthful checker, every accepted $\mathit{best}$ is a
genuine witness. Two corollaries pin down the conservative half:
\texttt{observe\_unknown\_preserves\_best} (an \texttt{Unknown} never advances
$\mathit{best}$) and \texttt{observe\_unknown\_eq\_doesNotTrigger} (an
\texttt{Unknown} is observationally identical to a non-reproduction). This is the
formal statement that a flake or timeout can never be promoted to the witness.

\subsection{The no-\texttt{sorry} policy}
The verified core admits no proof holes: zero \texttt{sorry}, zero
\texttt{admit}, no new \texttt{axiom}, and no forced \texttt{UInt64} evaluation
(theorems are kept structural over an opaque step function so the kernel never
deep-recurses over 64-bit numerals). The policy is enforced mechanically by a
token checker and by \texttt{lake build} with \texttt{warningAsError=true}, and
the whole Lean build is gated by CI --- a change that introduces a hole or breaks
a proof cannot merge.

% ===========================================================================
\section{Reference-leads differential}\label{sec:differential}
% ===========================================================================
Proving properties of a Lean \emph{port} is only useful if the port faithfully
models the executable Haskell. \rc\ closes that gap with a reference-leads
discipline: \textbf{the Haskell implementation is the authority.} A generator
(\texttt{scripts/gen-conformance.hs}) runs the real Haskell
\texttt{candidates}, \texttt{start}, and \texttt{observe} over a committed scope
of inputs --- the empty timeline, single faults of every intensity-bearing
constructor at interior and boundary (cap) values, \tool{abort} (no weaker form),
passthrough decisions, and multi-fault timelines --- and emits the results both
as a JSON golden and as Lean terms.

Two Lean theorems then re-check those terms \emph{inside the kernel} by
\texttt{native\_decide}: \texttt{candidates\_conforms\_on\_candidateScope}
(the Lean candidate set reproduces the reference's exactly) and
\texttt{runShrink\_conforms\_on\_trajectoryScope} (a fuel-bounded Lean driver
under a deterministic oracle shrinks to exactly the reference's converged
$\mathit{best}$). On the Haskell side, a golden-snapshot test re-derives the same
corpus and fails if the committed bytes do not reproduce. The abstract theorems
hold for \emph{all} timelines; the corpus pins the \emph{concrete} behavior to
the reference over the committed scope. We are explicit about the trust boundary:
\texttt{native\_decide} trusts Lean's native evaluator, and neither the corpora
nor shared constants prove end-to-end equivalence of the Haskell implementation,
the IO shell, and the Lean model.

This same discipline runs across all five layers of \rc's verified core --- the
SplitMix64 keystream, the fault scheduler, the output-tree equivalence oracle,
the minimizer, and the replay checks --- each shipping a faithful Lean port,
kernel-checked abstract theorems, and a \texttt{native\_decide} conformance
theorem against a reference-generated corpus.

% ===========================================================================
\section{Related work}\label{sec:related}
% ===========================================================================
The algorithmic lineage is covered in \S\ref{sec:background}. On the
\emph{verification} axis --- the one this paper claims novelty on --- the nearest
prior art is formally verified \emph{compiler} optimization. Compositional
Optimizations for CertiCoq~\cite{paraskevopoulou2021certicoq} gives a
Coq-mechanized, semantics-preserving correctness proof for optimization passes
(including a ``shrink reduction'' in the inlining sense). It is thematically
adjacent --- machine-checked termination and correctness of a program transform
--- but it is not a test-case reducer, and its ``shrinking'' is a compiler
optimization, not failure-preserving input minimization. We are not aware of a
widely-known machine-checked delta-debugging or test-case reducer carrying
termination and witness-preservation proofs, which is why we think the artifact
is worth reporting; we make no priority claim, and would welcome pointers to
prior verified reducers from reviewers.

Empirically, recent work cautions against assuming any one shrinking technique
dominates: an evaluation across multiple generator families and workloads
(ETNA~\cite{etna2026shrinking}) finds QuickCheck's structural shrinking usually
faster and competitive on quality, with generator family and workload dominating
outcomes --- a result we take seriously in \S\ref{sec:discussion}.

% ===========================================================================
\section{Discussion: which upgrade?}\label{sec:discussion}
% ===========================================================================
The \rc\ shrinker is a greedy, 1-minimal, value-level reducer. Three well-studied
directions could improve it; the constraint unique to a high-assurance tool is
that \emph{the verification artifact must survive the change}.

\paragraph{Iterate to a fixpoint (DDMIN*).}
Re-running the search from $\mathit{best}$ until a full pass accepts nothing
escapes the first local optimum for a strictly smaller 1-minimal
result~\cite{vince2022iterating}. This is nearly free and, critically,
\emph{preserves every existing proof}: each pass is another instance of the same
terminating, witness-preserving machine, and the outer loop descends the same
$\mu$ across passes. We recommend this first.

\paragraph{Probabilistic query selection (ProbDD).}
Reordering \texttt{pending} by expected information gain~\cite{wang2021probdd}
attacks \rc's weakest point --- query count against the flaky oracle its
\texttt{Unknown} handling already pays for --- and slots in without touching the
candidate set. But its speedup comes precisely from \emph{abandoning}
1-minimality, the property we went to the trouble of proving. For a tool whose
pitch is trustworthiness, that is a bad default. It belongs, if anywhere, as an
opt-in fast mode that explicitly drops the 1-minimality claim.

\paragraph{Choice-sequence reduction.}
Hypothesis-style internal shrinking~\cite{maciver2020hypothesis} is elegant and
its shortlex objective is a close cousin of our $\mu$. But ETNA undercuts the
prior that it is strictly better, and adopting it would relocate the
witness-preservation obligation from a value-level deterministic predicate to
``generator re-execution preserves the witness'' --- a materially harder Lean
proof. Not worth it unless timelines grow enough internal structure that flat
chunking wastes queries; today they do not.

\paragraph{Grammar-aware reduction.} Least relevant: it pays off on
tree-structured inputs~\cite{misherghi2006hdd,sun2018perses}, whereas a \rc\
timeline is a flat fault sequence where chunk deletion is already near-optimal.

% ===========================================================================
\section{Limitations}\label{sec:limitations}
% ===========================================================================
The proofs are about the pure core, not the running tool. The IO shell, the live
server's event scheduling, the Python drivers, and the external checker are
outside the Lean kernel's reach; a fixed seed does not fix a live server's event
trace. Witness preservation is conditional on a \emph{truthful} checker --- if
the user's checker reports \texttt{Triggers} on a candidate that does not in fact
reproduce the bug, the shrinker will faithfully preserve the wrong witness.
Deletion-1-minimality is a local optimum (\S\ref{sec:background}); we do not claim
global minimality. The \texttt{native\_decide} corpus checks trust Lean's native
evaluator, and correspondence between the Haskell reference and the Lean model is
established on finite corpora and sampled QuickCheck properties, not by a proof of
full equivalence.

% ===========================================================================
\section{Conclusion}\label{sec:conclusion}
% ===========================================================================
Shrinking is the step that makes a chaos finding actionable, and its two
load-bearing properties --- termination and witness preservation --- are
safety-critical to the claim ``this is the minimal cause.'' We showed that these
properties can be \emph{proved}, not merely engineered, for a practical shrinker
that combines structural deletion with a domain-specific intensity axis, and that
the proved model can be kept honest against the executable by a reference-leads
differential. The reduction algorithm is ordinary delta debugging; the
contribution is the verification artifact and the discipline that keeps it
faithful. We offer this draft for mixed human/AI peer review ahead of deeper
implementation work on the upgrade paths of \S\ref{sec:discussion}.

% ===========================================================================
{\footnotesize
\bibliographystyle{plainnat}
\bibliography{refs}
}

\end{document}
